\section{Importance Judgments Favor Stated Candidates}\label{sec:gap}

The importance prompt $u(c,x)$ asks whether $c$ is one of the most important concepts to remember for possible future questions and requests a single yes or no (App.~\ref{app:prompts}); Qwen3 uses thinking-disabled prompts. Let $Y,N$ contain distinct first tokens of the case and leading-space forms of \emph{yes} and \emph{no}, respectively:
\begin{equation}
\begin{split}
s_\theta(c;x)={}&\log\!\sum_{t\in Y}p_\theta(t\mid u(c,x))\\
&-\log\!\sum_{t\in N}p_\theta(t\mid u(c,x)),\\
V_\theta(c;x)={}&\sigma\!\left(s_\theta(c;x)\right).
\end{split}
\label{eq:score}
\end{equation}
The untrained model's score $V$ is the importance report; SDA (\S\ref{sec:method}) ranks by $s_\theta$ with its adapter enabled. App.~\ref{app:prompts} defines the 1--10 importance rating.

The workspace comprises verbalizable concepts the model is currently reasoning with \citep{anthropic2026workspace}. The workspace readout decodes each layer's residual state with the logit lens \citep{nostalgebraist2020logitlens,anthropic2026workspace}. Let $h_\ell(x)$ be the residual state at the final context position after layer $\ell$, of $L$ layers. Final normalization $\mathrm{LN}$ and unembedding $W_U$ give logits $z_\ell(x)=W_U\,\mathrm{LN}(h_\ell(x))$; token ranks are $\mathrm{rank}_\ell(t)=1+|\{v:z_{\ell,v}>z_{\ell,t}\}|$. The workspace readout is the peak reciprocal rank:
\begin{equation}
W(c;x)=\max_{\substack{1\leq\ell\leq L\\ t\in F(c)}}
\frac{1}{\mathrm{rank}_\ell(t)}.
\label{eq:wrr}
\end{equation}
$F(c)$ contains distinct first tokens of the candidate's bare and leading-space forms, each in original and capitalized form. One forward pass over $x$ supplies the required states. Output likelihood is $\mathrm{LL}(c;x)=\log p_{\mathcal M}(c\mid x)/|c|$, where $|c|$ is token length.

The importance report favors stated candidates when stated and unstated candidates compete. Among unstated candidates, its pooled AUC is 0.880--0.897 across six models, against 0.592--0.697 for the workspace readout (Table~\ref{tab:withinunstated}). Report AUC falls to 0.658--0.736 on the statedness-balanced pool and below chance, 0.322--0.427, on the original pool, where its statedness AUC is 0.582--0.688. On the statedness-balanced pool, correlations with statedness are $+0.90$ for single-pass ConsistencyGate, $+0.85$ for sampled ConsistencyGate and $+0.75$ for A-MAC utility.

The workspace readout ranks useful concepts higher than the importance report on the statedness-balanced pool (pooled AUC 0.719--0.812) and above chance on the original pool (six-model mean 0.563). The stated-bridge control, which appends ``The passage also mentions \emph{concept}.'' naming the useful concept, raises mean workspace readout AUC to 0.956; report AUC rises from 0.373 to 0.504 (Figure~\ref{fig:hero}a, Table~\ref{tab:statedcontrol}). Naming a distractor lowers both on Qwen3-8B (App.~\ref{app:naming}).

Rewording the importance prompt to say the passage may omit the useful concept changes the direction of this preference rather than removing it (App.~\ref{app:wording}).

The model's intermediate layers encode information about unstated useful concepts that its importance report does not express; SDA turns this information into admission decisions (\S\ref{sec:method}).
